% part: intuitionistic-logic
% chapter: introduction
% section: constructive-reasoning

\documentclass[../../../include/open-logic-section]{subfiles}

\begin{document}

\olfileid{int}{int}{cr}

\olsection{الاستدلال البنائي}

بخلاف توسيعات المنطق الكلاسيكي بالمؤثرات الموجهية أو كمّيات الرتبة
الثانية، يكون المنطق الحدسي «غير كلاسيكي» من حيث إنه يقيّد المنطق
الكلاسيكي. فالمنطق الكلاسيكي \emph{غير بنائي} بوجوه شتى. ويُراد
بالمنطق الحدسي أن يضبط نمطًا «أكثر بنائية» من الاستدلال يميز نوعًا من
الرياضيات البنائية. وقد توضح الأمثلة الآتية بعض الدوافع الكامنة وراءه.

افترض أن شخصًا ادعى أنه عيّن عددًا طبيعيًا~$n$ له الخاصية الآتية: إذا
كان $n$ زوجيًا فحدسية ريمان صحيحة، وإذا كان $n$ فرديًا فحدسية ريمان
خاطئة. يا لها من بشرى!{} فكون حدسية ريمان صحيحة أو غير صحيحة من
المسائل الكبرى المفتوحة في الرياضيات، ويبدو أن هذا الشخص قد رد
المسألة إلى مجرد حساب، أي إلى تعيين ما إذا كان عدد محدد زوجيًا أم لا.

فما القيمة السحرية لـ~$n$؟ يصفها كما يأتي: $n$ هو العدد الطبيعي الذي
يساوي $2$ إذا كانت حدسية ريمان صحيحة، ويساوي $3$ خلاف ذلك.

فتطلب غاضبًا استرداد مالك. من وجهة نظر كلاسيكية، يحدد الوصف السابق
فعلًا قيمة وحيدة لـ$n$؛ لكن ما تريده حقًا هو قيمة لـ$n$ معطاة
\emph{صراحة}.

ولنأخذ مثالًا آخر قد يكون أقل افتعالًا. انظر في السؤال الآتي. نعلم أنه
يمكن رفع عدد غير نسبي إلى قوة نسبية والحصول على نتيجة نسبية؛ فمثلًا
$\sqrt{2}^2 = 2$. أما الأقل وضوحًا فهو إمكان رفع عدد غير نسبي إلى قوة
\emph{غير نسبية} والحصول على نتيجة نسبية. تجيب المبرهنة الآتية
بالإيجاب:

\begin{thm}
يوجد عددان غير نسبيين $a$ و$b$ بحيث يكون $a^b$ نسبيًا.
\end{thm}

\begin{proof}
انظر في $\sqrt{2}^{\sqrt{2}}$. إذا كان هذا العدد نسبيًا فقد انتهينا:
نأخذ $a = b = \sqrt{2}$. وإلا فهو غير نسبي. وعندئذ
\[
(\sqrt{2}^{\sqrt{2}})^{\sqrt{2}} = \sqrt{2}^{\sqrt{2} \cdot
  \sqrt{2}} = \sqrt{2}^2 = 2,
\]
وهو عدد نسبي. ولذلك نأخذ في هذه الحالة
$a=\sqrt{2}^{\sqrt{2}}$ و$b=\sqrt 2$.
\end{proof}

هل يُعد هذا برهانًا صحيحًا؟ يرى معظم الرياضيين أنه كذلك. لكن فيه مرة
أخرى شيء غير مُرضٍ قليلًا: فقد برهنا وجود زوج من الأعداد الحقيقية له
خاصية معينة من دون أن نستطيع بيان \emph{أي} زوج هو. ويمكن برهنة
النتيجة نفسها بحيث يُعطى الزوج $a$، $b$ \emph{في} البرهان: خذ
$a = \sqrt{3}$ و$b = \log_3 4$. فالعدد~$\sqrt3$ غير نسبي، وكذلك
$\log_3 4$ غير نسبي؛ إذ لو كان $\log_3 4=m/n$ لعددين طبيعيين موجبين
$m,n$، لكان $3^m=4^n=2^{2n}$، وهذا يناقض وحدانية التحليل إلى عوامل
أولية. ومن ثم يكون كل من $a$ و$b$ غير نسبي، ولدينا
\[
a^b = \sqrt{3}^{\log_3 4} = 3^{1/2 \cdot \log_3 4} = (3^{\log_3
  4})^{1/2} = 4^{1/2}= 2,
\]
لأن $3^{\log_3 x} = x$.

صُمم المنطق الحدسي لضبط نمط من الاستدلال لا يسمح بخطوات كالخطوة
الواردة في البرهان الأول. فبرهنة وجود $x$ يحقق~$!A(x)$ تعني أن عليك
إعطاء~$x$ محدد، وإعطاء برهان على أنه يحقق $!A$، كما في البرهان الثاني.
وبرهنة أن $!A$ أو $!B$ تصدق تستلزم قدرتك على برهنة إحداهما.

وبلغة صورية، المنطق الحدسي هو ما نحصل عليه إذا قيدنا نظام
!!a{derivation} للمنطق الكلاسيكي بطريقة معينة. ومن وجهة النظر
الرياضية، ليست هذه إلا أنظمة استنباطية صورية؛ لكنها، كما ذكرنا، ترمي
إلى ضبط نمط من الاستدلال الرياضي. ويمكن اعتبار هذا النمط هو الاستدلال
الذي يبرره رأي فلسفي معين في الرياضيات (كحدسية براور)، أو اعتباره نمطًا
«أكثر عيانية» وإرضاء من الاستدلال الرياضي (على نهج بنائية بيشوب)، كما
يمكن مناقشة ما إذا كان الوصف الصوري يضبط الدافع غير الصوري أم لا.
ولكن مهما تكن مواقفنا الفلسفية، يمكننا دراسة المنطق الحدسي بوصفه
منطقًا معروضًا عرضًا صوريًا؛ ولأسباب شتى يجد كثير من المناطقة
الرياضيين متعة في دراسته.

\end{document}
